\subsection{Sesquilinear maps.}

In this section, unless expressly stated otherwise, we denote by A and B two rings, by E a left A-module and by F a left B-module; we denote by \(b \to b^{\mathrm{J}} (b \in \mathrm{B})\) an anti-automorphism of B, that is, a bijection from B onto itself satisfying \((b + c)^{\mathrm{J}} = b^{\mathrm{J}} + c^{\mathrm{J}}\) and \((bc)^{\mathrm{J}} = c^{\mathrm{J}}b^{\mathrm{J}}\) for all \(b, c\) in B; we write \(\mathbf{J}'\) instead of \(\mathbf{J}^{-1}\) . We denote by G a \((\mathbf{A}, \mathbf{B})\) -bimodule (no. 1).

DEFINITION 4. --- A map \(\Phi\) from E \(\times\) F into G is said to be right-sesquilinear for J if it satisfies conditions (1), (2), (3) (def. 1, no. 1) as well as

\begin{enumerate}
\def\labelenumi{(\arabic{enumi})}
\setcounter{enumi}{6}
\tightlist
\item
  \(\Phi(x, by) = \Phi(x, y). b^J\) for all \(x \in \mathbf{E}, y \in \mathbf{F}\) and \(b \in \mathbf{B}\) .
\end{enumerate}

If J is the identity (which requires that B be commutative), we recover the notion of a bilinear map.

Let \((e_i)_{i\in I}\) and \((f_k)_{k\in K}\) be two families of elements of E and F, and let \((a_i)_{i\in I}\) and \((b_k)_{k\in K}\) be elements of A and B that are zero except for a finite number of them. We then have

\[
\Phi (\sum_ {i} a _ {i} e _ {i}, \sum_ {k} b _ {k} f _ {k}) = \sum_ {i, k} a _ {i} \Phi (e _ {i}, f _ {k}) b _ {k} ^ {\mathrm{J}}.\tag{8}
\]

As in the case of a bilinear map, the elements \(\Phi(e_i, f_k)\) determine \(\Phi\) in a unique way when \((e_i)\) and \((f_k)\) are systems of generators, and can be chosen arbitrarily when \((e_i)\) and \((f_k)\) are bases of E and F; when \((e_i)\) and \((f_k)\) are finite bases, we say that \((\Phi(e_i, f_k))\) is the matrix of \(\Phi\) with respect to these bases.

As for bilinear maps, one defines on the set of right sesquilinear maps (for J) from E × F into G a bimodule structure over the centers of A and B. One defines the notion of inverse image of a sesquilinear map by the same formula as for a bilinear map. We shall moreover see that the study of sesquilinear maps can be reduced to that of bilinear maps.

DEFINITION 5. --- Let B be a ring, F a left (resp. right) B-module and J an antiautomorphism of B. One denotes by F \(^{j}\) the right (resp. left) B-module having the same underlying additive group as F and in which the external composition law is \((b, y) \to b^{J'}y\) (resp. \((b, y) \to yb^{J'}\) ) \((b \in B, y \in F, J' = J^{-1})\) .

With the notations of Definition 5, a linear map from \(\mathbf{F}^{\mathrm{j}}\) into a right (resp. left) B-module H is therefore identified with a Z-linear map \(u\) of F into H satisfying

\[
u (b y) = u (y) b ^ {J} \quad (\text { resp. } u (y b) = b ^ {J} u (y)) \quad (b \in B, y \in F).
\]

The map \(u\) from F into H is a semilinear map of F into H relative to J (chap. II, App. I, no. 1), if one regards J as an isomorphism of the opposite ring B \(^{0}\) of B onto B, and F as a B \(^{0}\) -module on the right (resp. on the left).

Similarly, a right-sesquilinear map \(\Phi\) (for J) from E × F into G, where F is a left B-module, identifies with a bilinear map from E × F \(^{J}\) into G; if the latter is degenerate on the right (resp. degenerate on the left, nondegenerate), one says that \(\Phi\) is degenerate on the right (resp. degenerate on the left, nondegenerate).

Remark. --- Let A and B be two rings, \(J_{1}\) an anti-automorphism of A, M a right A-module, N a right B-module, and G an (A, B)-bimodule. One says that a map \(\Phi\) from M \(\times\) N into G is left-sesquilinear for J \(_{1}\) if it is Z-bilinear and if it satisfies

\[
\Phi (x a, y b) = a ^ {\mathrm{J} _ {1}} \Phi (x, y) b \quad (x \in \mathrm{M}, y \in \mathrm{N}, a \in \mathrm{A}, b \in \mathrm{B}).\tag{9}
\]

Such a map is identified with a bilinear map from \(M^{J_{1}} \times N\) into G. We will often leave it to the reader to transpose to left sesquilinear maps the definitions and properties given for right sesquilinear maps; when we speak of a sesquilinear map (without specifying), it will be a right sesquilinear map.

